\documentclass[11pt]{article}
\usepackage[a4paper,margin=2.3cm]{geometry}
\usepackage{microtype}
\usepackage{enumitem}
\usepackage{xcolor}
\usepackage{titlesec}
\usepackage{tabularx}
\usepackage{booktabs}
\usepackage[most]{tcolorbox}
\usepackage[hidelinks]{hyperref}
\usepackage{parskip}

\definecolor{accent}{HTML}{000000}
\definecolor{soft}{HTML}{000000}
\definecolor{boxbg}{HTML}{FFFFFF}
\definecolor{boxline}{HTML}{000000}

\titleformat{\section}{\large\bfseries\color{accent}}{\thesection.}{0.5em}{}
\titleformat{\subsection}{\normalsize\bfseries\color{accent}}{}{0em}{}
\titlespacing*{\section}{0pt}{1.7ex}{0.7ex}

\setlist[itemize]{leftmargin=1.4em,topsep=2pt,itemsep=3pt}
\setlist[enumerate]{leftmargin=1.7em,topsep=2pt,itemsep=3pt}

% Plain-language definition box
\newtcolorbox{plain}[1]{
  colback=boxbg, colframe=boxline, boxrule=0.5pt, arc=2pt,
  left=8pt,right=8pt,top=5pt,bottom=5pt, breakable,
  title={\small\bfseries In plain terms: #1}, fonttitle=\color{black},
  coltitle=black, colbacktitle=boxbg, attach title to upper=\hspace{0pt}}

\newcommand{\reportnote}[1]{%
  \par\vspace{3pt}%
  {\small\color{soft}$\hookrightarrow$\ \textbf{Where this is in the report:} #1}\par\vspace{2pt}}

\hypersetup{pdftitle={Thesis Brief for Supervisor}}

\begin{document}
\pagestyle{empty}

\begin{center}
  {\LARGE\bfseries\color{accent} What Our Research Did --- A Brief for My Supervisor}\\[5pt]
  {\large Privacy-Preserving, Fairness-Aware Federated Fraud Detection on Hyperledger Fabric}\\[7pt]
  {\small\itshape This document tells the story of the project in plain language. Technical terms are
  explained the first time they appear, and again in the glossary at the end. Every section says where
  the same material lives in the full thesis report. All numbers come from re-running our actual code.}
\end{center}
\vspace{0.3em}
\hrule
\vspace{0.7em}

%==================================================================
\section{The problem we started with}

Banks want to catch fraud. The trouble is that any single bank only sees \emph{its own} customers'
transactions. A fraud pattern that is spread across many banks --- the same stolen card tried at ten
different institutions --- is almost invisible to each bank on its own. If the banks could combine
their data, their fraud detectors would get dramatically better.

But they can't simply pool the data. Privacy laws forbid sharing raw customer records; banks are
commercial rivals who don't want to hand over their data; and there is no neutral, trusted party that
everyone is willing to let hold the combined dataset. This is the deadlock our project tackles.

\begin{plain}{the deadlock}
Everyone would benefit from pooling data, but nobody is allowed to, wants to, or trusts anyone enough
to actually do it. We need the \emph{benefit} of pooling without the \emph{act} of pooling.
\end{plain}

There is a known idea that promises a way out, called \textbf{federated learning}.

\begin{plain}{federated learning (FL)}
Instead of sending your data to a central place, each bank trains a model on its own data, in-house,
and only shares the trained model's \emph{numbers} (its ``weights'') --- never the raw transactions.
A coordinator combines everyone's numbers into one shared, smarter model. The data never leaves the
building; only lessons learned from it do.
\end{plain}

Federated learning sounds like it solves everything. But the moment you try to actually run it between
banks that distrust each other, four new problems appear that the textbook glosses over:

\begin{enumerate}
  \item \textbf{Who runs the coordinator, and why trust them?} Whoever combines the models holds a
  position of power. Rivals won't accept one bank (or a third party) being in charge.
  \item \textbf{Why contribute good data?} If everyone receives the same final model, a lazy bank can
  contribute garbage (or nothing) and still get the benefit --- a ``free-rider.'' Good contributors
  need to be \emph{measured} and \emph{paid} fairly.
  \item \textbf{The shared numbers still leak.} It turns out a model's weights can, with effort, be
  reverse-engineered to reveal things about the data they were trained on. So even the numbers need
  protection.
  \item \textbf{Someone might sabotage the shared model} on purpose by sending poisoned numbers.
\end{enumerate}

The tempting pitch is that you can fix all four at once by bolting on four well-known tools ---
a \emph{blockchain} for trust, \emph{differential privacy} for leakage, \emph{Shapley values} for fair
payment, and \emph{Krum} for sabotage-resistance (all explained below). The result, in theory, is a
system that is private, fair, secure, and decentralised all at the same time.

\textbf{Our research question is simply: does that pretty promise actually hold up when you build it for
real?} We took an existing, published fraud detector (the DB-BOA-ADTCN model of Prabanand \& Thanabal,
2025) and moved it into this real federated, blockchain setting to find out.

\reportnote{Chapter~1 (Introduction) sets up this problem and lists our objectives; Chapter~2
(Literature Review) explains where each of the four tools comes from and what gap we fill.}

%==================================================================
\section{The four tools, in plain language}

Before the story continues, here is what each bolt-on tool actually does. These four names recur
throughout the report, so it is worth fixing them in mind.

\begin{plain}{blockchain / Hyperledger Fabric}
A shared digital ledger that no single party controls --- every participant keeps a copy, and entries
can't be secretly changed. It removes the need for a trusted middleman. \textbf{Hyperledger Fabric} is
a specific, business-grade blockchain platform; we run a real one, not a pretend one. The small program
that defines our rules on it (logging results, handing out reward tokens) is called a ``chaincode.''
\end{plain}

\begin{plain}{differential privacy (DP)}
A mathematical way to add carefully-calibrated random ``noise'' (small fudge factors) to the numbers a
bank shares, so that no individual customer's data can be backed out of them --- while the overall
model is still useful. There is a dial, written $\epsilon$ (epsilon): \emph{small $\epsilon$ = lots of
noise = strong privacy but blurrier numbers; large $\epsilon$ = little noise = weak privacy but clearer
numbers.} The whole trade-off in this thesis lives on this dial.
\end{plain}

\begin{plain}{Shapley values}
A fair way, borrowed from game theory, to answer ``how much did each player actually contribute to the
team's success?'' We use it to score how much each bank's data improved the shared model, so rewards
can be split fairly. We represent each round's scores as a short list of numbers, $\varphi$ (phi).
\end{plain}

\begin{plain}{Krum}
A defence against sabotage. When combining everyone's model numbers, instead of naïvely averaging them
(where one wildly poisoned set can drag the average off), Krum picks the one set of numbers that is
most ``in agreement'' with the majority and throws out the outlier. It has a mathematical guarantee ---
but only when there are enough honest participants relative to attackers.
\end{plain}

\reportnote{Chapter~4 (Proposed Methodology) describes how we wired these four together into one
pipeline, and Chapter~3 states our ``threat model'' (exactly which attacks we defend against).}

%==================================================================
\section{What we built}

We didn't just describe this on paper --- we implemented the whole stack and made it run end to end:

\begin{itemize}
  \item \textbf{The detector.} A fraud-classifier model (ADTCN) whose settings are tuned automatically
  by a search algorithm called \textbf{DB-BOA}, so a human doesn't have to hand-tune it. The same
  algorithm also picks which node ``leads'' each blockchain round.
  \item \textbf{The federation.} Three simulated banks --- BankA, BankB, BankC --- each train locally,
  and their models are combined through the pipeline: \emph{add privacy noise (DP) $\rightarrow$ reject
  sabotage (Krum) $\rightarrow$ score contributions (Shapley)}.
  \item \textbf{The rewards.} Each bank's Shapley contribution score automatically determines its cut of
  a pool of reward ``tokens,'' recorded on the blockchain. Contribute more, earn more --- no human
  judgement involved.
  \item \textbf{The blockchain.} A real Hyperledger Fabric network runs underneath, logging every round
  and keeping the token accounts.
\end{itemize}

But the \emph{point} of the project was never just ``assemble the stack.'' The real question was the
hard one from Section~1: \textbf{when you actually combine privacy, fairness, and security, do they
cooperate --- or do they get in each other's way?}

\reportnote{Chapter~4 (Proposed Methodology) --- the full pipeline, the DB-BOA tuning, and the chaincode.}

%==================================================================
\section{The story of what happened (the heart of the thesis)}

Here is the turn that makes this a piece of research rather than just an engineering demo.

When we turned on \emph{strong} privacy --- the privacy dial set to a tight $\epsilon = 1.0$ (lots of
protective noise) --- and ran the complete system honestly, it didn't just get a bit worse. It
\textbf{fell apart}. The combined model became useless: it stopped distinguishing fraud from non-fraud
at all, effectively just labelling everything the same way.

\begin{plain}{why it fell apart}
The privacy noise was being sprinkled over a \emph{huge} set of numbers (the model's ${\sim}111{,}000$
weights). At strong privacy, that much noise drowned the real signal entirely. And Krum, the sabotage
defence, couldn't rescue it --- by design Krum keeps just \emph{one} bank's noisy numbers and discards
the rest, which throws away the averaging that would otherwise cancel some of the noise out.
\end{plain}

A weaker project, faced with this, would have quietly turned the privacy dial down until the demo looked
good again, and never mentioned it. We did the opposite. We decided this failure \emph{was} the finding:
\textbf{privacy and fairness are fundamentally in tension, and the interesting scientific question is to
map exactly where, and why, and whether it can be fixed properly.}

So we asked: is privacy itself the problem, or is it \emph{where we were applying the privacy noise}? We
had been noising the giant pile of model weights and then trying to read each bank's contribution
through that fog. What if, instead, we left the model weights alone and applied the privacy noise
\emph{directly to the small contribution scores} ($\varphi$) --- just three numbers instead of a
hundred thousand?

That one change of \emph{where} the noise goes is the key result of the thesis. Protecting three numbers
needs far less noise than protecting a hundred thousand, so fair rewards survive at a much stronger
privacy setting.

\begin{plain}{the fix, in one line}
Don't blur the whole model to hide a contribution --- just blur the contribution score itself. Same
privacy promise, a tiny fraction of the damage.
\end{plain}

\reportnote{Chapter~5 (Result Analysis), \S6.2 (the collapse) and \S6.6 (the privacy-vs-fairness
characterisation and our fix). The reframed contribution is also stated in Chapter~1.}

%==================================================================
\section{What we found --- the actual numbers}

Every figure here comes from re-running our code (saved in \texttt{final\_report\_data/}).

\subsection{The detector works (and we report it honestly)}
On the public ULB credit-card fraud dataset, our detector scored \textbf{99.85\% accuracy}. But accuracy
is misleading here, because only 0.17\% of transactions are fraud --- a lazy model that says ``never
fraud'' already scores 99.83\%. So we report \textbf{MCC} instead.

\begin{plain}{MCC (Matthews Correlation Coefficient)}
A single honest score from $-1$ to $+1$ that only looks good if the model gets \emph{both} fraud and
non-fraud right. $0$ means ``no better than guessing.'' It can't be fooled by the ``always say normal''
trick. Our detector scores \textbf{MCC 0.677} --- genuinely useful.
\end{plain}

Honest caveat we state openly: our \emph{auto-tuned} detector (0.677) is actually a touch \emph{worse}
than a hand-set default (0.785). So DB-BOA's real value is \emph{removing the manual tuning work}, not
beating a human expert --- and we say exactly that.

\subsection{The centrepiece result: the privacy-vs-fairness fix}
Moving the privacy noise from the model weights onto the small contribution scores improves the privacy
dial at which rewards stay \emph{fair} (in the right order) by about \textbf{60 times}:

\begin{center}
\small
\begin{tabularx}{0.92\linewidth}{Xcc}
\toprule
& \textbf{Old way (noise on weights)} & \textbf{Our way (noise on scores)} \\
\midrule
Privacy needed for fair rewards & $\epsilon^* = 3000$ (basically no privacy) & $\epsilon^* = 50$ (usable privacy) \\
Reward fairness at $\epsilon=50$ (corr., $-1$ to $1$) & $\approx -0.29$ (rewards backwards!) & $\approx +0.95$ (rewards correct) \\
\bottomrule
\end{tabularx}
\end{center}

In words: the old approach only gave fair rewards if you switched privacy off almost entirely
($\epsilon=3000$); our approach gives fair rewards at a genuinely private setting ($\epsilon=50$). The
``$-0.29$'' is striking --- at that setting the old method actually pays the \emph{worst} contributors
the \emph{most}. Our method gets it right ($+0.95$, almost perfect ordering).

\subsection{Security holds where it's supposed to}
We tested Krum against four kinds of sabotage attack, in the conditions where its mathematical guarantee
is meant to apply (enough honest banks). Krum \textbf{caught and rejected the attacker in all 8 test
cases}, keeping the shared model at \textbf{$\approx$99.9\%} balanced accuracy. Without Krum, the worst
attack dragged a plain average down to \textbf{$\approx$87.5\%}. So the security claim is real --- within
stated limits.

\subsection{Each ingredient's effect, side by side (an honest ablation)}
\begin{center}
\small
\begin{tabularx}{0.95\linewidth}{Xccl}
\toprule
\textbf{Setup} & \textbf{Accuracy} & \textbf{MCC} & \textbf{What it tells us} \\
\midrule
Plain averaging (FedAvg) & 99.72\% & 0.569 & decent baseline \\
+ Krum (anti-sabotage) & 99.92\% & \textbf{0.776} & best single setup \\
+ strong privacy ($\epsilon{=}1.0$) & 99.83\% & 0.000 & privacy alone collapses it \\
Full system at $\epsilon{=}1.0$ & --- & ${\approx}0$ & the collapse from Section~4 \\
\bottomrule
\end{tabularx}
\end{center}
We \emph{show} the collapse rather than hiding it --- it is the evidence for our central point.

\begin{plain}{FedAvg}
The standard, naïve way to combine models in federated learning: just average everyone's numbers
together. Simple, but it has no defence against sabotage or privacy leakage on its own --- which is why
the other tools get added.
\end{plain}

\subsection{The blockchain numbers are measured, not invented}
We measured our real Fabric network rather than guessing: each agreement round takes about
\textbf{2.1 seconds} (a built-in timing floor of the platform), and it sustains about \textbf{40
transactions per second} before hitting a known congestion limit. (Earlier drafts had quoted made-up
figures; these replace them.)

\reportnote{Chapter~5 (Result Analysis) --- §6.1 detector, §6.2 ablation, §6.6 the privacy-vs-fairness
fix, the Byzantine-robustness section (Krum 8/8), and the Fabric-measurement section.}

%==================================================================
\section{What we can honestly claim as our contribution}

Three things, strongest first, in words that will survive a tough examiner:

\begin{enumerate}
  \item \textbf{A new privacy-preserving reward mechanism} --- the idea of putting the privacy noise on
  the \emph{contribution score} instead of the \emph{model weights}, which makes fair rewards survive
  ${\sim}60\times$ more privacy. \emph{This is the one genuinely new piece.}
  \item \textbf{A clear characterisation of the privacy-vs-fairness trade-off} --- real evidence that
  naïvely combining strong privacy with fair-reward scoring makes rewards \emph{worse than random}, plus
  a map of exactly where that breakdown happens. \emph{A legitimate ``negative result'' finding.}
  \item \textbf{A complete, working, honestly-measured system} --- all the pieces integrated and running
  on a real blockchain. \emph{This is integration of known tools, and we label it as such.}
\end{enumerate}

We are careful \emph{not} to overclaim: we did not invent differential privacy, Krum, Shapley values, or
Fabric, and we do \emph{not} claim to be ``first'' to put fair rewards on a blockchain --- an earlier
system (FedCoin, 2020) did that, and we build on it.

\begin{plain}{why ``negative results'' still count as research}
Showing clearly and with evidence that a popular combination \emph{doesn't} work the way people assume
--- and explaining why --- is a real scientific contribution. It saves the next team from the same trap
and points to the right fix. Honesty here is a strength, not a weakness.
\end{plain}

\reportnote{Chapter~1 (contribution list) and Chapter~6 (Conclusion).}

%==================================================================
\section{Where we fall short --- our limitations}

We state these plainly in the report. Examiners trust honest limitations far more than a flawless-looking
story.

\begin{itemize}
  \item \textbf{The ``three banks'' are really one dataset split three ways.} We didn't have three real
  banks' data, so we split one public dataset into three portions. That means we never truly tested banks
  whose data looks genuinely \emph{different} from each other.
  \item \textbf{The fairness/reward defence only works in some situations.} It can isolate a coordinated
  \emph{majority} of cheaters, but it does \emph{not} catch a lone attacker or a quiet free-rider (who
  contributes nothing but isn't actively malicious). Fixing free-riding needs an extra rule we leave for
  future work.
  \item \textbf{``Secure'' and ``Scalable'' are bounded claims, not blanket ones.} Krum's guarantee only
  holds when honest participants sufficiently outnumber attackers; a very subtle attacker that stays close
  to honest behaviour could still slip through. And ``scalable'' refers to our \emph{contribution-scoring}
  method, not the blockchain's raw speed (which has a clear ceiling).
  \item \textbf{Some results are from a simulation, not a live deployment.} The blockchain speed numbers
  and the leader-picking behaviour come from a single-computer test network, not a real multi-bank,
  multi-city deployment. The \emph{patterns} transfer; the exact numbers wouldn't.
  \item \textbf{A small measurement caveat on the contribution scores.} They're computed on held-out data
  drawn from the same pool as the training data, so there's a minor, disclosed circularity --- not an
  independent bank's data.
\end{itemize}

\reportnote{Chapter~5 / Chapter~6 --- the ``Threat model'' paragraph and the ``Limitations and
disclosures'' block (drafted in \texttt{final\_report\_data/REWRITE\_08\_limitations\_disclosures.md}).}

%==================================================================
\section{If my supervisor reads only one paragraph}

\begin{quote}
\itshape
We built a system that lets rival banks jointly train a fraud detector without sharing raw data, on a
real blockchain, with privacy protection, fair contribution-based rewards, and sabotage resistance. Our
central finding is that these good properties \textbf{fight each other}: turn privacy up to a serious
level and the fair-reward mechanism breaks so badly it pays the worst contributors the most, and the
model collapses. We pinned down exactly why --- the privacy noise was being applied in the wrong place
--- and fixed it by noising the small contribution score instead of the huge model, which restores fair
rewards at roughly $60\times$ stronger privacy. The sabotage defence works as promised within stated
limits, and our blockchain numbers are measured rather than invented. So the contribution is: one new
mechanism, one honest characterisation of a real trade-off, and a complete working system --- with its
limitations (single dataset, conditional reward defence, simulated network) stated openly.
\end{quote}

%==================================================================
\section{The whole story in one sitting (a plain walkthrough)}

If the earlier sections still feel like a pile of separate parts, here is the same research told as a
single, continuous story --- no jargon assumed.

\subsection{The setting}
Imagine several banks that each want to catch fraudulent transactions. A better fraud detector needs more
data --- but \textbf{no bank will hand its customers' records to the others} (privacy law, and they are
rivals). The standard way out is \textbf{federated learning}: each bank trains a model on its
\emph{own} data and shares only the trained model, never the raw data. A coordinator combines these into
one stronger shared model. Everyone benefits; nobody exposes raw data.

\subsection{Three things you need at once --- and they fight}
A real federated system between distrustful banks has to do three jobs at the same time, and they pull
against each other:
\begin{enumerate}
  \item \textbf{Fair pay (incentive).} One bank has excellent data, another has almost none. Pay them
  equally and the good one walks away. So you must \emph{measure} how much each bank actually helped
  (\textbf{Shapley value}) and pay it accordingly, recorded tamper-proof on a \textbf{blockchain}.
  \item \textbf{Privacy.} Even sharing the \emph{model} can leak information about individuals. The
  defence is \textbf{differential privacy}: add random \emph{noise}. More noise = more privacy.
  \item \textbf{Security.} A malicious bank might send a poisoned model to wreck the shared one. The
  defence is \textbf{Krum}, which discards the suspicious outlier instead of blindly averaging.
\end{enumerate}

\subsection{The conflict at the heart of the thesis}
Here is the tension everything hangs on: \textbf{privacy sabotages fairness.} To get privacy you add
noise --- but that same noise scrambles the contribution measurement (Shapley). When the measurement is
scrambled, the blockchain \textbf{pays the wrong banks}: the free-rider gets rewarded, the hard worker
gets shortchanged. The incentive becomes a lie. So the honest research question becomes:
\emph{``When does this reward system actually stay honest, instead of paying the wrong people?''}

\begin{plain}{why this is the novelty}
All the parts already exist (federated learning, Shapley, privacy, Krum, blockchain). What is new is
\emph{mapping out exactly when the combination breaks --- and fixing it.} That is a characterisation
contribution, not a new-algorithm claim.
\end{plain}

\subsection{What we discovered (the negative result)}
We measured it. Adding the privacy noise the ``obvious'' way --- onto the \textbf{model weights}
($\sim$100{,}000 numbers) --- the rewards only stay correct when privacy is \emph{almost switched off}
($\epsilon\approx3000$). In plain terms: with the obvious design you can have correct rewards \emph{or}
real privacy, but not both. That is a real, useful finding on its own.

\subsection{What we built (the fix --- the centrepiece)}
Then the key insight: the noise does so much damage because it is dumped onto 100{,}000 numbers. But you
do not need to privatise the weights to protect the reward calculation --- you only need to privatise the
\textbf{contribution scores themselves}, which are just a handful of numbers (one per bank). So we flip
the order: compute the honest scores on the \emph{clean} models first, \emph{then} add the privacy noise
to those few numbers before publishing them. Same privacy guarantee, but noising a handful of numbers
instead of 100{,}000 --- so the damage drops by roughly two orders of magnitude. Honest rewards now
survive at a genuinely private setting ($\epsilon\approx50$).

\begin{plain}{the punchline}
\emph{Where} you inject the privacy noise relocates the honesty boundary by about $60\times$. We turned a
``you can't have both'' negative result into a ``here's how you get both'' positive one.
\end{plain}

\subsection{Why the other title words are earned}
\textbf{Secure} --- Krum catches poisoning that plain averaging misses. \textbf{Scalable} --- the exact
Shapley sum is too expensive, so we use a sampling shortcut that scales, plus we \emph{measured} real
blockchain throughput. \textbf{Consensus / RL} --- the blockchain integration, with a learning agent
choosing which node leads each round. These are honest integration and measurement, not new inventions.

\reportnote{This walkthrough condenses Chapters~1, 4 and~5; the fix and its numbers are §6.6.}

%==================================================================
\section{Glossary (quick reference)}

\begin{description}[leftmargin=0pt,style=nextline,itemsep=3pt]
  \item[Federated learning (FL)] Training a shared model from many parties' data without any of them
  sending their raw data --- only model numbers are shared.
  \item[Hyperledger Fabric] A real, business-grade blockchain platform; we run one as the trusted,
  no-middleman coordinator. The on-chain rule program is the ``chaincode.''
  \item[Differential privacy (DP) / $\epsilon$] Adding tuned random noise to shared numbers so no
  individual can be identified. $\epsilon$ is the dial: small = strong privacy + blurry; large = weak
  privacy + clear.
  \item[Shapley value ($\varphi$)] A fair score of how much each participant contributed to the team
  result; we use it to split rewards.
  \item[Krum] A defence that picks the most-agreed-with model and discards sabotage outliers, with a
  guarantee when honest participants are numerous enough.
  \item[FedAvg] The plain baseline: just average everyone's model numbers, with no built-in defences.
  \item[MCC] An honest accuracy score ($-1$ to $+1$) that isn't fooled when fraud is rare; $0$ = guessing.
  \item[ADTCN / DB-BOA] The fraud-detector model, and the search algorithm that auto-tunes it (and picks
  blockchain leaders), respectively.
  \item[Byzantine attack] Jargon for a participant deliberately sending bad/poisoned data to wreck the
  shared model.
  \item[Free-rider] A participant who takes the shared model's benefit while contributing little or
  nothing useful.
\end{description}

\vspace{0.5em}
\hrule
\vspace{0.4em}
{\footnotesize\color{soft}\itshape Full report chapter map: Ch1 Introduction $\cdot$ Ch2 Literature
Review $\cdot$ Ch3 Requirements, Impacts \& Constraints $\cdot$ Ch4 Proposed Methodology $\cdot$ Ch5
Result Analysis $\cdot$ Ch6 Conclusion. Verified numbers: \texttt{final\_report\_data/}. Anticipated
defence questions: \texttt{defense\_questions.md}.}

\end{document}
